\section{Introduction}
\label{sec:intro}

Regulated organisations increasingly want to train and adapt foundation models on confidential data that
cannot leave a jurisdiction, using whatever accelerator capacity they own across a few metropolitan sites.
Our running example is a Canadian financial institution with GPU sites in Toronto, Montr\'eal and a U.S.\
region, confidential (here: synthetic) financial documents, live video/audio assistance, and checkpoints
that must stay in Canada.

In such deployments, three subsystems are usually tuned by three different teams. The ML platform team
chooses the \emph{placement} $\pi$ (which ranks run at which site, and in what ring order). The network team
chooses the \emph{routing} $r$ over a small number of WAN paths, typically with a traffic-engineering (TE)
system in the style of B4~\citep{b4} or SWAN~\citep{swan} that sees demand matrices but not jobs. The
storage team chooses where checkpoints go ($s$), subject to residency and geo-redundancy policies. Each
team optimises its own objective under assumptions about the other two.

This paper asks one question:
\begin{quote}\itshape
Does cross-layer optimisation of compute placement, network routing and checkpoint storage significantly
reduce training time, step-time variance and failure-recovery time compared with independently tuned
subsystems, under reproducible workload and failure conditions?
\end{quote}

We answer it with a reproducible emulation study rather than a production deployment. A three-site WAN with
two disjoint inter-site planes runs real FRRouting eBGP/iBGP/OSPF control planes and real kernel forwarding
with \texttt{tc} shaping. Data-parallel training ranks run in per-site network namespaces, so every Gloo
collective and every checkpoint byte crosses the emulated WAN. A controller steers traffic only through BGP
policy (local-preference per destination prefix), and it can re-place the job through
checkpoint-and-restart migration. We compare a joint (π, r, s) planner against static routing, ECMP, an
\emph{independently tuned} configuration (best placement for the nominal network plus a B4-style
min-max-utilisation TE optimiser that sees measured demand), and an oracle fed ground-truth network state.
The comparison covers seven failure types plus a steady state, with seeded background traffic,
interleaved method order, and five seeds per cell.

\paragraph{Contributions.}
\begin{itemize}
\item A calibrated, analytical step-time model for synchronous data-parallel training over a
multi-path WAN. It couples ring order, per-edge path choice, background load and checkpoint flushes, and an
exact joint planner over $(\pi, r, s)$ for small site counts (\S\ref{sec:design}).
\item An open, scripted laboratory: netns + FRR WAN, failure injection, telemetry, and a job manager with
migration (\S\ref{sec:impl}). Every number in this paper is regenerated by \texttt{make reproduce-full}.
\item An evaluation of joint versus independent control across failure types, with ablations of each
telemetry signal and of migration (\S\ref{sec:eval}). We report negative and mixed results as found.
\item Supporting subsystems, evaluated separately: a crash-consistent, content-addressed checkpoint
engine; an mTLS/SPIFFE control plane with live certificate rotation and authenticated, encrypted
checkpoints; TLA+ models of the checkpoint and rotation protocols; and small extensions for media, kernels,
post-training and privacy (\S\ref{sec:ext}).
\end{itemize}

\paragraph{What this study is not.} It does not use GPUs or real WAN links. It is a
CPU-only testbed with scaled-down link capacities, chosen so that communication is a large share of
step time, as it is for geo-distributed training. Absolute times do not transfer to production; we claim
only relative effects inside the regime we calibrate and report (\S\ref{sec:limits}).
